\documentclass[12pt]{article}
\usepackage[utf8]{inputenc}
\usepackage[T1]{fontenc}
\usepackage{hyperref}
\title{Zhuang, Liujiang Language Resources for K-12 Curriculum Development}
\author{Compiled by Claude Sonnet 5.5 (claude-sonnet-5-5) for Andr\'as Kornai}
\date{}
\begin{document}
\maketitle

\section{Introduction}
Liujiang Zhuang (ISO 639-3 zlj; Glottocode liuj1238; Chinese name \emph{Liujiang Zhuangyu}) is a Tai-Kadai language of Guangxi, China, and belongs to the Zhuang macrolanguage. Glottolog places it in the Lianshan-Liujiang subgroup with Lianshan Zhuang and rates it as not endangered. The searches found only catalogue entries. No Liujiang-specific textbooks or programmes were found. Standard Zhuang, a different variety, has its own educational materials which were not searched here.

\section{Primary Educational Resources}
No primers, textbooks or school programmes specific to Liujiang Zhuang were found.

\section{Linguistic and Reference Materials}
\begin{itemize}
\item Glottolog record for ISO code zlj. \url{https://glottolog.org/resource/languoid/iso/zlj}
\item Detect Language entry for Liujiang Zhuang. \url{https://www.detectlanguage.com/languages/zlj}
\item Wikipedia, ``Zuojiang Zhuang languages'': related Zhuang varieties (Zuojiang is not Liujiang). \url{https://en.wikipedia.org/wiki/Zuojiang_Zhuang_languages}
\end{itemize}

\section{Religious and Community Resources}
\begin{itemize}
\item Global Recordings Network entry for Liujiang Zhuang; it listed no recordings at the time of the search. \url{https://www.globalrecordings.net/en/language/zlj}
\end{itemize}

\section{Digital and Audio Resources}
\begin{itemize}
\item Global Recordings Network, ``Zhuang, Liuzhou Liucheng'': a related entry that has recordings and is also tagged zlj. Its exact relation to Liujiang Zhuang was not verified. \url{https://globalrecordings.net/en/language/6129}
\end{itemize}

\section{Curriculum Development Context}
With about 1.6 million speakers, the language has a large community, but no Liujiang-specific teaching materials were found. A curriculum would probably draw on standard Zhuang materials and adapt them to local Liujiang speech, with local fieldwork needed.
\end{document}